\documentclass[11pt,a4paper]{article}
\usepackage[margin=2.2cm]{geometry}
\usepackage[T1]{fontenc}
\usepackage[utf8]{inputenc}
\usepackage{lmodern}
\usepackage{microtype}
\usepackage{amsmath,amssymb}
\usepackage{booktabs}
\usepackage{array}
\usepackage{enumitem}
\usepackage{caption}
\usepackage{subcaption}
\usepackage{graphicx}
\usepackage{xcolor}
\usepackage{pgfplots}
\pgfplotsset{compat=1.17}
\usepackage[hidelinks]{hyperref}

\definecolor{cours}{HTML}{3B6EA8}
\definecolor{coccam}{HTML}{3D9A6B}
\definecolor{cfill}{HTML}{B5532A}
\definecolor{cgrey}{HTML}{777777}
\newcommand{\hit}{\textcolor{coccam}{$\bullet$}}
\newcommand{\miss}{\textcolor{cgrey}{$\circ$}}
\newcommand{\code}[1]{\texttt{\small #1}}
\makeatletter\newcommand{\inputtable}[1]{\@@input #1 }\makeatother

\title{\textbf{Region-level teachers in our pipeline}\\[2pt]
\large \code{--clip-mode langsplat} and the skip-vs-fill test for uncovered pixels}
\author{Radiance FYP --- semantic embedding, branch \code{semantic/new}\\ \small intellisense08 (RTX 2080, 8\,GB), 16 October 2026}
\date{}

\begin{document}
\maketitle

\begin{abstract}
\noindent Occam's LGS run on LangSplat's features was the best design measured (25/29 objects, 1/6 absent objects firing,
14 Oct). This report covers two things. First, those region-level features are now a teacher mode of our own pipeline
(\code{--clip-mode langsplat}), lifted with our all-pixel lift, so they can be built, cached, queried and served like
every other table. Second, the open question of how to treat pixels no SAM mask covers is tested: \emph{skip} them, as
LangSplat and Occam do, or \emph{fill} them with the scale-averaged crop feature, as our B2 variant did. With skip, our
pipeline reaches \textbf{24/29 with 1/6 false positives} (dev set 5/5). That is within one query of Occam's chain, the
only difference being a cluster offset on one object. \textbf{Skip vs fill differs by one query} (24 vs 23 at the default
floor, 24 vs 24 at equal false positives), so filling does not explain why B2 reached only 21. The default lift mode is
unchanged (A); switching it is a decision for Suhan.
\end{abstract}

\section{What was built}
\begin{description}[leftmargin=1.2em,style=nextline,itemsep=2pt]
  \item[Teacher step, \code{--clip-mode langsplat} (cache tag \code{4f4bdb1a08})]
    \code{teachers.extract\_langsplat} runs LangSplat's preprocessing for every frame as a subprocess in the
    \code{occamlgs} env, which has LangSplat's SAM fork. That means SAM ViT-H masks at levels default / s / m / l, mask
    NMS, and one OpenCLIP ViT-B/16 embedding per masked crop, with LangSplat's settings, which are part of the cache tag.
    It resumes per frame. The 720 training-view feature files from the 13--14 Oct run were moved into the cache
    (symlinks remain at the old path); only the 30 + 49 evaluation-split frames were extracted (37 + 31\\,min). DINOv2
    grids are shared with the existing teacher cache (same model and width), so L-CD stays available.
  \item[Lift, \code{--uncovered skip|fill} (\code{lift.lift\_segment\_views})]
    Our gsplat lift, run per level $\ell\in\{s,m,l\}$. Each pixel takes the feature of the level-$\ell$ mask covering
    it; seg maps are sampled at the render's pixel centres. With $m_\ell(p)=1$ where a mask covers $p$:
    \begin{align*}
      \textit{skip:}\quad f_i^{(\ell)} &= \frac{\sum_{v,p} w_i(v,p)\,m_\ell(p)\,F_\ell(p)}{\sum_{v,p} w_i(v,p)\,m_\ell(p)}
        &&\text{(LangSplat / Occam's choice)}\\
      \textit{fill:}\quad f_i^{(\ell)} &= \frac{\sum_{v,p} w_i(v,p)\,\big[m_\ell F_\ell + (1-m_\ell) F_{\text{pyr}}\big](p)}{\sum_{v,p} w_i(v,p)}
        &&\text{(our B2's choice)}
    \end{align*}
    $F_{\text{pyr}}$ is the scale-averaged crop pyramid (teacher \code{d7a30d7264}). The skip denominator is one extra
    one-channel pass per level.
  \item[Table and query] Unchanged. The three levels are the table's three scale groups, so the query picks the level
    with the highest peak, as LangSplat's evaluation does. \code{--uncovered} is a variant key, so existing lift markers
    stay valid. The lift default remains \code{scales} (A).
  \item[Tests] Four new unit tests: skip and fill agree when every pixel is covered; skip never takes the uncovered side;
    fill mixes it in; seg-map sampling. Semantics suite 89 passed, Galley 116 passed.
\end{description}

\noindent Build commands (tables \code{lift\_ls} and \code{lift\_lsfill}):
{\footnotesize
\begin{verbatim}
python figs/semantic_pipeline.py --scene backroom --backend lift --clip-mode langsplat \
       --table-suffix _ls --uncovered skip
python figs/semantic_pipeline.py --scene backroom --backend lift --clip-mode langsplat \
       --table-suffix _lsfill --uncovered fill
\end{verbatim}
}

\section{Results}
\begin{table}[h]
\centering\small
\caption{Phase 5 sets (29 positives, 6 absent objects), query settings unchanged since Phase 1. Totals, then per scene:
hits, median localisation error (m), AUROC of peak relevancy.}
\label{tab:summary}
\resizebox{\linewidth}{!}{%
\begin{tabular}{llcc ccc ccc}
\toprule
& & & & \multicolumn{3}{c}{backroom} & \multicolumn{3}{c}{flightroom}\\
\cmidrule(lr){5-7}\cmidrule(lr){8-10}
Configuration & mode & hits & FP & hits & err & AUROC & hits & err & AUROC\\
\midrule
\inputtable{tables/summary.tex}
\bottomrule
\end{tabular}}
\end{table}

\begin{table}[h]
\centering\small
\caption{The only queries (L-C) on which these configurations differ: hit (\hit) or miss (\miss) and peak relevancy. b =
backroom, f = flightroom. All other 25 queries have the same outcome in all five.}
\label{tab:perq}
\begin{tabular}{l ccccc}
\toprule
Query & Ours A & Ours B2 & Occam full & LangSplat, fill & LangSplat, skip\\
\midrule
\inputtable{tables/per_query_diff.tex}
\midrule
\inputtable{tables/per_query_total.tex}
\bottomrule
\end{tabular}
\end{table}

\begin{table}[h]
\centering\small
\caption{Floor sweep (0.46--0.60, post-hoc on the evaluation sets): best hits with at most $k$ absent objects firing,
and hits at the default floor 0.55.}
\label{tab:env}
\begin{tabular}{lcccccc}
\toprule
$k$ absent objects with a candidate & 1 & 2 & 3 & 4 & 5 & at 0.55\\
\midrule
\inputtable{tables/envelope.tex}
\bottomrule
\end{tabular}
\end{table}

\begin{figure}[h]
\centering
\begin{tikzpicture}
\begin{axis}[width=0.8\linewidth, height=5.4cm, xlabel={relevancy floor}, ylabel={top-1 hits (of 29)}, ymin=14, ymax=27,
  legend style={font=\scriptsize, at={(0.5,-0.25)}, anchor=north, legend columns=3}, grid=major, grid style={gray!20}]
  \addplot[cours, thick, dashed, mark=square*, mark size=1.2pt] table[x=floor, y=A_h, col sep=comma] {data/sweep.csv};
  \addplot[cgrey, very thick, mark=diamond*, mark size=1.4pt] table[x=floor, y=Occamfull_h, col sep=comma] {data/sweep.csv};
  \addplot[coccam, very thick, mark=*, mark size=1.4pt] table[x=floor, y=LSskip_h, col sep=comma] {data/sweep.csv};
  \addplot[coccam, thick, dashed, mark=o, mark size=1.4pt] table[x=floor, y=LSskipflat_h, col sep=comma] {data/sweep.csv};
  \addplot[cfill, thick, mark=triangle*, mark size=1.4pt] table[x=floor, y=LSfill_h, col sep=comma] {data/sweep.csv};
  \legend{Ours A, Occam's full chain, LangSplat skip, {LangSplat skip, no level choice}, LangSplat fill}
\end{axis}
\end{tikzpicture}
\caption{Hits against the relevancy floor, both scenes combined.}
\label{fig:sweep}
\end{figure}

\section{Findings}
\begin{enumerate}[leftmargin=*,itemsep=3pt]
  \item \textbf{The region teachers carry over to our pipeline.} With our all-pixel lift and skip, 24/29 with one absent
        object firing (only backroom's ceiling fan), AUROC 0.92 / 1.00, dev set 5/5. Like Occam's chain, it recovers
        \emph{swivel chair}, \emph{camera tripod} and \emph{hose reel}, which A misses. The one difference from Occam's
        chain is flightroom's \emph{water jug}. Both find the jug's neighbourhood, but our cluster centroid lands outside
        the hit radius ("bad position"). This $\pm1$ between the all-pixel and centre-pixel lifts matches the 13 Oct
        controlled comparison.
  \item \textbf{Skip vs fill: one query.} The only query that differs is \emph{camera tripod}: peak 0.588 with skip,
        0.576 with fill. Filling the 6--17\,\% of uncovered pixels with crop features pulls a thin object's features
        towards its background and costs this marginal hit. At equal false positives both reach 24/29. The flat
        variants (no level choice) reach 25 with skip at one false positive and 24 with fill. Skip is the better choice,
        but the effect is small.
  \item \textbf{So B2's shortfall is not the fill.} B2 reached 21/29 and misses all three region-level gains. With fill
        costing at most one query, the gap must come from B2's other choices: SAM ViT-B without a crop layer and with SAM's stricter
        default thresholds (predicted IoU 0.86 vs 0.7), levels by mask area instead of SAM's own granularities, and
        no LangSplat mask NMS. All three can now be tested as teacher variants
        of this mode.
  \item \textbf{Level choice is not consistently useful.} With Occam's lift, choosing the level per query added two
        (25 vs 23 at one false positive). With our lift, the flat average is as good or better (25 vs 24). On 29 queries
        this is noise. The level choice should be kept configurable, and the default decided on a larger query set.
  \item \textbf{Cost.} The lift takes 18 / 21\,min per scene (three 512-channel levels plus DINO, about the same as A's
        lift). Occam's centre-pixel lift does the same in 30\,s. Feature extraction is the expensive part: about 5\,h per
        scene on the RTX 2080 (70\,s per 1080p image), done once and cached. Tables: 2.5 / 1.8\,GB, as A.
\end{enumerate}

\section{Decisions for Suhan and next steps}
\begin{enumerate}[leftmargin=*,itemsep=2pt]
  \item \textbf{Default lift mode.} \code{langsplat} with skip is the strongest mode in our own pipeline (24/29, 1/6;
        A: 22/29, 2/6), and its features are already cached for both scenes. Switching the default is one line
        (\code{default\_clip\_mode}). The cost is about 5 GPU-hours of extraction per new scene on the 2080.
  \item \textbf{Speed the lift} by adding Occam's centre-pixel estimator as an option (30\,s instead of 20\,min; the
        13 Oct comparison found no accuracy loss).
  \item \textbf{Close the B2 gap by ablation:} SAM ViT-B vs ViT-H, crop layer on/off, area levels vs SAM levels, NMS
        on/off. Each needs a new SAM extraction, so they are best run on the RTX 5060 Ti or overnight.
  \item \textbf{Galley:} \code{langsplat} tables load and query like any multi-scale table; the build should be exposed
        in Galley's semantics job form once the default is settled.
\end{enumerate}

\section{Files}
Raw results: \code{../raw/} (Phase 5 and development-set evaluator JSONs, lift statistics including per-level pixel
coverage, floor sweep, run log). Derived CSVs: \code{../csv/}. Tables: \code{tables/}, written by
\code{../scripts/collect.py}. Code: branch \code{semantic/new} (\code{lift.py}, \code{teachers.py},
\code{semantic\_pipeline.py}, \code{baselines/langsplat\_features.py}, tests).

\end{document}
